Maria was elk uur
in de vrije natuur

toen ze nog heel klein was.
Alsof het haar genas

van het eenzaam voelen.
Ze zocht paddenstoelen,

mos en wilde bloemen,
zonder ze te noemen,

maar ze op te nemen
of ze mee te nemen.

Ze kon struiken rossen.
Ze zwierf door de bossen

en sloop door de togel.
Als paradijsvogel

bracht ze zichzelf geluk.
Huma kwam in haar stuk

en combineerde een mens
met zijn eenwordingswens.

Deze holistische,
leerfilosofische

en zeer open aanleg
ging uiteindelijk weg.

“Is echt het heelal één?
Ik voel me heel alleen!”

Vader probeerde weer
Maria keer op keer

aan de kerk te slijten.
Tijdens zijn verwijten,

vroeg een priester toen bot
aan hem: “uit welke grot

ben jij ooit gekropen?”
De grond misgelopen,

dacht kleine Maria
dat elke paria

letterlijk in een grot,
steeds kon worden gespot.

Na haar afgeluister,
gelokt door het duister,

kroop ze in een spelonk,
die rottend benauwd stonk

en vond daar een bandiet,
die haar in leven liet.

Ze kroop ook in een grot
en ontdekte haar lot

middels een kluizenaar.
De man vertelde haar:

“Lief kind, de diepste schacht
herbergt de grootste pracht.

Verbloem nooit wat je denkt.
Want wat je aandacht schenkt,

juist door in te tomen,
zal je overkomen

en verklaar je geldig.
Je bent pas geweldig

als je de moed opbrengt
en alles overbrengt

in jouw boek van leven
en hebt opgeschreven!

Doch het op papier laat
en niets ermee aangaat.

Zo ontneem je het kwaad
zijn verwoestende daad.”

Maria kon zoeken
tussen vele boeken

naar gewenste kennis.
Het mocht van de abdis.

Ze leerde zo alert
dat ze er gek van werd.

Die epistemische
aandacht, socratische

bewustwording als wel
theoretisch gekwel,

maakte haar studiereis
vast geen leerparadijs.

Het klooster werd een hel.
Maar dat kende ze wel.

“Het leven uitzingen
brengt veel uitdagingen.

Als je dat accepteert
heb je leven geleerd!

Dan levert de leerhel
het menszijn zielsherstel,

door ermee om te gaan
en het zelf af te staan.

Dan, als de tijd verstrijkt,
wordt het goede verrijkt.”

Maria ging toen staan.
“Novice doe er wat aan

blijf er niet mee zitten!”
Ze ging erbij zitten.

Emmy stond in het licht
van Maria’s gezicht.

Ze voelde een raakvlak.
Alsof ze haar aansprak.

“Omdat onwetendheid
tot een lijdensweg leidt

is ‘leren’ bevrijden
van nodeloos lijden.

Een tranendal in bloei.”
Spirituele groei

via een ‘lesje leren’
kon Maria deren.

“Hoe kan men zelf leren
zonder te bekeren?

In vrijheid God kennen
zonder te verkennen?

Vraagt dat niet vertrouwen?
Dat men steeds kan bouwen

op genoeg veiligheid?
Verantwoordelijkheid

vraagt toch autonomie?”
Het verbaasde Emmy

dat Maria dit zei,
onvrij in de abdij.

Weer kwam een naïeve.
Maria haatte ze.

Die zuivere vreugde.
Een mens die toch deugde.

Ze sprak de novice aan,
maar wilde liever slaan.

Het gesprek liep anders
dan slechts wat uitbranders.

Maria werd in woord
in haar lijden gehoord.

De open debutant
bleek plots een geestverwant,

die haar kon aanhoren.
Het kon haar bekoren,

dus schaamde ze zich rot.
“Je maakt alles kapot!”,

zei een stem in haar hoofd,
die ze stil had geloofd.

Emmy luisterde mee
en ging er van de twee

grondig mee aan de slag.
Waardoor Emmy inzag

dat de stem kon baten.
Ze kreeg in de gaten

dat de stem vertelde
dat ze zichzelf kwelde.

En dat gaf haar de kracht
om haar ontspoorde macht

ten goede te voeren.
Ze wilde doorvoeren

en langzaam afdwingen
breinveranderingen.

Al Emmy’s constructen
uit haar brein mislukten

om Maria’s bangheid
te zien als haar bloeistrijd.

Emmy kwam tot inkeer.
“Ik weet niets zeker meer.

Maria wist het niet.
Hoe God haar ooit verliet

en in angst dompelde.
Ze daarna strompelde

door haar leven geleefd
hardleers en onbeleefd.

Ik zie weinig notie
in mijn incarnatie,

maar waar ben ik hierin?
Ben ik zelf de heldin

die wantrouwen omarmt.
Die zijn angsten verwarmt?”

Zo sprak de illusie:
“Ik zeg je mijn Emmy

ik ben je aan het testen
niet dwaas aan het pesten.

Falen, beter falen,
op je best gaan falen

is je afleerproces.
Geen doelen is de les.

Je leert perfectie af.
Falen is loon en straf.

Je vrees voelt als een hel.
Dit is echter een spel

om te kunnen leren.
Te manifesteren,

niet uit behoeftigheid,
maar uit verbondenheid

middels ieders lijden
en daaruit bevrijden.”

De dichter bij de brei
is dichterbij het ei

dat Columbus legde.
Het in rijm gezegde

is duaal gecodeerd,
zodat men beter leert.

Vandaar dat Maya dicht
omwille van inzicht.

Waar de donkere nacht
van de ziel Emmy bracht

zal Maya vertellen
in vele leerhellen,

waar Mara zal dollen
in meerdere rollen.
